\chapter{A Chapter Title 1}
\chapterstyle{deposit}

\section{Introduction}
There are over 14,000 school districts that operate in the United States today charged with governing American public schools. Over 13,000 of these are run independently of any other governmental entity's authority by locally appointed or elected officials.\footnote{Data from several sources including: \citep{RefWorks:325,RefWorks:321,RefWorks:322}}  These local officials make policy decisions in school auditoriums at sparsely attended evening meetings that affect everything from local property tax rates to graduation requirements to the location and type of schools to open and close. In a majority of communities they are the single largest employer. They are responsible for annual expenditures roughly equivalent to the budget of the US Department of Defense.\footnote{Data comes from the Common Core of Data on school district finances. School districts reported revenue over \$550 billion in fiscal year 2007, while the respective Department of Defense budget for 2007 was just over \$500 billion.} These decisions ripple through nearly every American community annually and have wide reaching implications for everything from property values to economic development and most directly, for the quality of education provided to the children within each community. In short, school boards make decisions that matter--especially for the nearly 50 million students that attend the schools governed by such boards.\footnote{2011 data from NCES. http://nces.ed.gov/fastfacts/display.asp?id=372} 

What are school boards? School boards are locally elected (or appointed) special purpose governments charged with managing public schools. They typically have five to nine members that are locally elected in at large, non-partisan, rarely contested elections \citep{RefWorks:321}. Boards are charged with managing the school district, raising revenue through local taxes, and setting the policy direction of local schools. Most boards hire a professional manager--a district superintendent of schools most often--to handle the administration of the district and to advise the board on relevant questions of policy. This relationship between the board and the superintendent is a critical feature of local education governance.

School boards present a tremendous opportunity for political scientists to explore a number of puzzles about American democracy. School boards are just one example of the tens of thousands of local governments that exist for both general and special purposes and that make decisions that affect communities across America. Boards can shed light on the implications of different electoral rules on representativeness in a way that cannot be explored in House districts where single-member districts are virtually set in stone \citep{RefWorks:199,RefWorks:216,RefWorks:238}. Political scientists can evaluate the efficiency of centralizing control of a policy area under a strong executive like a mayor in lieu of a special government body like a school board--a common governance reform for public schools \citep{RefWorks:323, RefWorks:211}. By examining school boards, political scientists can even evaluate things like the response style of politicians with much larger sample sizes and more variation than can be found in any legislative body in the country \citep{RefWorks:313}. Furthermore, school boards are the last stop of the web of overlapping funding, authority, and regulations in education policy--providing an opportunity to evaluate the way these federal, state, and local entanglements play out under a variety of conditions \citep{RefWorks:318}. Yet, school boards are rarely used as a unit of analysis for such subjects, despite the staggering number of school boards and elected school board members, and the variety therein in the United States today.

Despite the wide reaching scope and importance of their decisions, local school boards remain woefully understudied by social sciences and particularly by political scientists. A search of JSTOR's political science section for articles with the term ``school governance'' over the last decade (2000-2010) results in five articles on American school boards, and forty-five articles on international school governance. That an overwhelming majority of school boards are elected, and that as a group elected school board officials make up one of the largest groups of elected officials in the American democratic system, makes this dearth of research all the more problematic. 

\section{The Puzzle}
Although there is wide room for exploration of many political phenomena through the study of school boards, this dissertation will focus on a single critical piece of the puzzle--how does a sudden, contentious, and deeply partisan state level education policy shock affect local governments? Specifically:
\vspace{-6mm}
\begin{enumerate}
\begin{singlespace}
\setlength\itemsep{-6pt}
\item How does such a policy shock affect the level of participation in school board elections both in terms of voter turnout and in terms of challengers to incumbent board members?
\item How does that shock change school board policymaking and what are the predictors of that change in the board?
\item Does sudden polarization and partisanship of state level politics lead to evidence of partisanship on school boards?
\end{singlespace}
\end{enumerate}
\vspace{-3mm}
Answering these questions can provide great insight into educational governance more broadly and specifically the interaction of school boards with other levels of government. The recent political events in Wisconsin have created a natural experiment within the state that allows causal leverage on these questions for the first time. By comparing the behavior of candidates, board members, boards, and voters before and after the policy shock it is possible to evaluate the causal impact of state level policy on the politics of school boards. At the same time, much insight can be gained into school boards within a state by looking at the universe of school boards and the dynamics of school board elections across an extended period of time \citep{RefWorks:192}. Finally, by tying board policy, electoral activity, and public policy preferences together along a single issue dimension identified exogenously to all school boards a spatial model (of ideological space) of school board elections can be evaluated empirically for the first time. 


\vspace{.1in}
\section{Contribution of Findings}

This dissertation contributes to the understanding of American politics in several ways. First, and foremost, it provides much needed focus on one of the most common democratic institutions in the American political system--school boards. 

Second, it tests theories of voter turnout and candidate participation in diverse contexts far from the traditional venues of state and federal legislative offices. Do traditional political science theories generalize to other elected offices, or are they confined to explaining the unique conditions of state and federal legislative and executive races? Do the traditional theories of school board politics explain the response of local political entities to policy change at the state or federal level?

This dissertation will help policymakers understand the pressures on and the governing capacity of local school boards in making policy to improve student achievement. Understanding the correlates of board policy changes along specific issue dimensions--employee compensation plans and responses to deep budget cuts--is an important step toward further understanding of the role of school districts in making and carrying out education policy in relation to state and federal policy. It also allows researchers to start exploring the role of stability in leadership on a diverse array of school district outcomes ranging from reform attempts, to employee satisfaction, to ultimately student outcomes.  



